%!TEX root = ./../../springer-mono-format/book.tex
%!TEX root = ./../../thesis-trier.tex

\IN{monotonic automaton}
An automaton $\semiautomat{A}{Q}{\Sigma}{\delta}$
is called \emph{monotonic} if there exists a linear order~$\le$ on the states
such that $p \le q$ implies $\applytrans{\delta}{p}{x} \le \applytrans{\delta}{q}{x}$
for all $p, q \in Q$ and $x \in \Sigma$.
See \cref{constr_sync:fig:monotonic_aut} for two example monotonic automata.

Such an order is generally not unique.  
For example, if all transitions are self-loops in an $n$-state automaton, $n!$ orderings are possible, 
and for each state, there exists an ordering in which it is either the minimal or the maximal element

Suppose $q = \delta(q, u^i)$ for $q \in Q$, $i \ge 1$ and $u \in\Sigma^*$ in a monotonic $\semiautomat{A}{Q}{\Sigma}{\delta}$.
We have $q \le \delta(q, u)$ or $\delta(q, u) \le q$.
In the first case, we get $q \le \delta(q, u) \le \delta(q, u^2) \le \ldots \le \delta(q, u^{i-1}) \le q$,
so that $q = \delta(q, u)$.
A similar argument yields $q = \delta(q, u)$
for $\delta(q, u) \le q$.
So, no word can permute a subset of states
non-trivially.\footnote{Automata with this property are called counter-free or aperiodic~\cite{McNaughton71}.}
In particular, no permutation automata is monotonic.

\begin{figure}[htb]
	\centering
	\begin{subfigure}[b]{0.45\textwidth}
		\centering
			\begin{tikzpicture} [node distance = 5.5em, on grid, auto, ->, >=triangle 45]
			
			\node (q0) [state, inner sep=5pt,minimum size=0pt] {$1$};
			\node (q1) [state, right of=q0,inner sep=5pt,minimum size=0pt] {$2$};
			\node (q2) [state, right of=q1,inner sep=5pt,minimum size=0pt] {$3$};
			\node (q3) [state, right of=q2,inner sep=5pt,minimum size=0pt] {$4$};
			
			\path[->] (q0) edge [bend left] node {$a$} (q1)
			(q1) edge [bend left] node {$b$} (q0)
			(q0) edge [loop above] node {$b$} (q0)
			(q2) edge [bend left] node {$b$} (q1)
			(q1) edge [bend left] node {$a$} (q2)
			(q2) edge [bend left] node {$a$} (q3)
			(q3) edge [bend left] node {$b$} (q2)
			(q3) edge [loop above] node {$a$} (q3);
		\end{tikzpicture}
		\caption{Linear order on states is $1 < 2 < 3 < 4$}
	\end{subfigure}
	\hfill
	\begin{subfigure}[b]{0.45\textwidth}
		\centering
	\begin{tikzpicture} [node distance = 5.5em, on grid, auto, ->, >=triangle 45]
		
		\node (q0) [state, inner sep=5pt,minimum size=0pt] {$1$};
		\node (q1) [state, right of=q0,inner sep=5pt,minimum size=0pt] {$2$};
		\node (q2) [state, right of=q1,inner sep=5pt,minimum size=0pt] {$3$};
		
		\path[->] (q0) edge [bend left] node {$a$} (q1)
		(q1) edge [bend left] node {$b$} (q0)
		(q0) edge [loop above] node {$b$} (q0)
		(q1) edge node {$a$} (q2)
		(q2) edge [loop above] node {$a,b$} (q2);
	\end{tikzpicture}
		\caption{Linear order on states is $1 < 2 < 3$} \label{prelim:fig:monotonic}
	\end{subfigure}
	\caption[Two synchronizing monotonic automata]{Two synchronizing monotonic automata.}
	\label{constr_sync:fig:monotonic_aut}
\end{figure}

With respect to synchronizability Ananichev and Volkov~\cite{DBLP:journals/tcs/AnanichevV04}
have established the following tight bound (tightness is 
shown by the automata from Figure~\ref{prelim:fig:unary_tight}, 
which are monotonic).

\begin{proposition}[Ananichev and Volkov~\cite{DBLP:journals/tcs/AnanichevV04}]
	\label{prelim:prop:monotonic_rank_r}
	Let $\semiautomat{A}{Q}{\Sigma}{\delta}$ be an $n$-state monotonic
	automaton. If there exists a word of rank $r$ in $\mathcal A$,
	then there exists a word of length
	at most $n - r$ whose rank is at most $r$
	and this bound is tight.
\end{proposition}